\chapter{Dynamics: equations, modes, and trustworthy computation}
\label{learn:dynamics}

\begin{learningcard}{Prerequisites and guiding experiment}
Use derivatives and integrals from Chapter~\ref{learn:calculus}, eigenvectors
from Chapter~\ref{learn:linear_algebra}, and conservation from
Chapter~\ref{learn:multivariable}. Imagine a warm bar cooling toward a
uniform temperature. We need an evolution equation, initial and boundary
data, a way to solve a controlled example, and a check that a computer's
approximation respects the mechanism.
\end{learningcard}

\section{An equation of motion needs initial data}
An ordinary differential equation (ODE) relates an unknown function of one
variable to its derivatives. For $y'=f(t,y)$, an initial condition
$y(t_0)=y_0$ chooses a trajectory when an existence and uniqueness theorem
applies. Continuity in $t$ and local Lipschitz continuity in $y$ near the
initial point provide a standard local guarantee. Lipschitz continuity means
$|f(t,u)-f(t,v)|\leq L|u-v|$ locally for a fixed finite constant $L$.
The guarantee is local: solutions can grow without bound in finite time,
as $y'=y^2$, $y(0)=1$, with $y=1/(1-t)$ illustrates.

\section{Worked example: exponential relaxation}
Suppose excess temperature $y$ obeys $y'=-ky$, where $k>0$ has units of
inverse time. Multiplying by $e^{kt}$ gives
$(e^{kt}y)'=e^{kt}(y'+ky)=0$. Therefore
\[
 y(t)=y_0e^{-kt}.
\]
This integrating-factor argument includes $y_0=0$, a case division by $y$
would have excluded temporarily. The time to reduce a positive excess by
half is $t_{1/2}=\log2/k$. With $k=0.2$ per minute and $y_0=10$ degrees,
the excess after five minutes is $10/e$ degrees, about $3.679$.
The exponential law is a model assumption. Its success in one regime does
not establish the same coefficient for every material or temperature range.

An equilibrium of an autonomous equation $y'=f(y)$ is a value $y_*$ with
$f(y_*)=0$. It is stable when sufficiently small initial perturbations stay
small for all forward time for which the comparison is made; it is
asymptotically stable when nearby trajectories also approach it. In the
relaxation model, the exact factor $e^{-kt}$ proves both properties at zero.

\section{Worked example: linear stability has directional content}
For a constant matrix system $z'=Az$, an eigenvector $v$ with eigenvalue
$\lambda$ produces the solution $z(t)=ce^{\lambda t}v$. If a basis of
eigenvectors exists, expand the initial vector in that basis and evolve
each coefficient. Consider
\[
 A=\begin{pmatrix}-2&1\\1&-2\end{pmatrix}.
\]
The vectors $(1,1)^T$ and $(1,-1)^T$ have eigenvalues $-1$ and $-3$.
For $z(0)=(2,0)^T$, the decomposition gives
\[
 z(t)=e^{-t}(1,1)^T+e^{-3t}(1,-1)^T.
\]
Both modes decay, and the slower symmetric mode dominates at late times.
A matrix with an eigenvalue of positive real part instead has a growing
mode. In finite dimensions, every eigenvalue having negative real part
ensures asymptotic stability even when generalized eigenvectors are needed.
Eigenvalues on the imaginary axis require more information.

For a smooth nonlinear system, write $z=z_*+\eta$ near an equilibrium.
Differentiability gives $\eta'=Df(z_*)\eta+o(\|\eta\|)$.
Strictly negative real parts imply local asymptotic stability under the
standard linearization theorem. A zero eigenvalue makes that test
inconclusive: $y'=-y^3$ attracts nearby trajectories, while $y'=y^3$
repels them, and both derivatives at zero vanish.

\begin{learningcard}{Historical situation: heat and a research program}
Fourier's \emph{Th\'eorie analytique de la chaleur} (1822) opens with a
preliminary discourse placing heat among mathematical investigations of
natural phenomena. The text develops equations and trigonometric expansions
for thermal problems. Its introductory account also describes the submission
of work to the Institut. That setting makes scientific institutions part of
the documented story of circulation and assessment.\footnote{Joseph Fourier,
\emph{The Analytical Theory of Heat}, translated by Alexander Freeman
(1878), Preliminary Discourse and Chapter I; University of Michigan
Historical Mathematics Collection:
\url{https://quod.lib.umich.edu/u/umhistmath/abr1558.0001.001}.}
The source documents a program and its presentation. Modern convergence and
solution theorems specify which Fourier manipulations are justified.
\end{learningcard}

\section{From a conservation law to the heat equation}
A partial differential equation (PDE) involves an unknown function of several
variables and its partial derivatives. In one space dimension let $u(x,t)$
be temperature, $c>0$ constant volumetric heat capacity, and $J$ heat flux.
Local energy conservation without sources says $cu_t+J_x=0$. Fourier's
constitutive law $J=-K u_x$ assumes a constant conductivity $K>0$ and flux
opposite the temperature gradient. Combining the two gives
\[
 u_t=\kappa u_{xx},\qquad\kappa=K/c>0.
\]
Here $\kappa$ has units of length squared per time. Conservation supplies
a balance; the constitutive law supplies material behavior. They are separate
assumptions that experiments can examine.

On $0\leq x\leq L$, impose zero endpoint temperature and initial profile
$u(x,0)=\sin(\pi x/L)$. Direct differentiation verifies the exact solution
\[
 u(x,t)=e^{-\kappa(\pi/L)^2t}\sin(\pi x/L).
\]
More generally, a finite linear combination of sine modes evolves by
multiplying mode $n$ by $e^{-\kappa(n\pi/L)^2t}$. Higher spatial frequencies
decay faster. Infinite expansions require convergence and regularity
conditions; an arbitrary interchange of a sum and a derivative needs proof.
For sufficiently smooth solutions with zero endpoint temperatures,
integration by parts gives
\[
 \frac{d}{dt}\frac12\int_0^L u^2\,dx
 =\kappa\int_0^Luu_{xx}\,dx=-\kappa\int_0^L u_x^2\,dx\leq0.
\]
The boundary term vanishes because $u$ vanishes there. This squared-norm
quantity supplies a stability check; total heat may leave through the ends.

\section{Stable dynamics and unstable time steps}
Forward Euler replaces $y'$ by a forward difference. For $y'=-ky$,
\[
 y_{m+1}=(1-k\Delta t)y_m.
\]
The exact solution decays for every positive time step, but the numerical
factor has magnitude at most one only when $0\leq k\Delta t\leq2$.
Strict decay requires $0<k\Delta t<2$; the endpoint two gives undamped
sign alternation. Nonnegativity for positive initial data requires the stronger
condition $k\Delta t\leq1$. At equality the iterate is zero; strict positivity requires $k\Delta t<1$.
Accuracy, stability, and sign preservation are distinct.

With equally spaced spatial points and time levels, the centered-space,
forward-time heat scheme is
\[
 u_j^{m+1}=u_j^m+r(u_{j+1}^m-2u_j^m+u_{j-1}^m),\qquad
 r=\frac{\kappa\Delta t}{\Delta x^2}.
\]
For periodic data, insert a mode $e^{ij\theta}$, where $i^2=-1$.
Its amplification factor is
$G(\theta)=1-4r\sin^2(\theta/2)$. The condition $0\leq r\leq1/2$
guarantees $|G|\leq1$ for every frequency; frequencies near $\pi$ expose
failure beyond that bound. Spatial refinement by a factor two therefore
requires a time step at most one quarter as large under this sufficient
stability rule.

A computation should record initial data, boundary rules, coefficients,
mesh widths, and time steps. Compare a known mode with its exact solution,
repeat on refined grids, and measure conservation under a boundary condition
that actually conserves the quantity. These checks test the discretization
and implementation; establishing a physical model additionally needs
observations and uncertainty estimates.

\section{Exercises with full solutions}
\paragraph{1. Test two Euler time steps.}
For $y'=-2y$, compare $\Delta t=0.4$ and $\Delta t=1.1$.

\paragraph{Solution.}
The amplification factors are $1-2(0.4)=0.2$ and
$1-2(1.1)=-1.2$. The first gives $y_m=0.2^my_0$, decaying with
preserved sign. The second gives $y_m=(-1.2)^my_0$, alternating and
growing in magnitude for nonzero initial data. The exact solution
$y_0e^{-2t}$ decays in both comparisons, identifying the growth as numerical.

\paragraph{2. Follow two heat modes.}
On $[0,\pi]$ with zero endpoints and $\kappa=1$, evolve
$u(x,0)=\sin x+2\sin3x$.

\paragraph{Solution.}
Each mode solves the equation separately, and linearity gives
$u(x,t)=e^{-t}\sin x+2e^{-9t}\sin3x$. Substituting $t=0$ verifies
the initial condition; both sines vanish at the endpoints. Two spatial
derivatives multiply the modes by $-1$ and $-9$, matching the time
derivatives. The coefficient ratio is $2e^{-8t}$, so the higher-frequency
contribution decays faster.

\paragraph{3. Prove discrete mass conservation.}
For the periodic heat scheme with $N$ points, prove preservation of
$M^m=\Delta x\sum_{j=0}^{N-1}u_j^m$ in exact arithmetic.

\paragraph{Solution.}
Summing the update gives
$M^{m+1}-M^m=r\Delta x\sum_j(u_{j+1}^m-2u_j^m+u_{j-1}^m)$.
Periodic indexing makes each shifted sum equal to $\sum_j u_j^m$,
so the difference is zero. Floating-point addition can introduce rounding
drift, which a program should measure. Conservation alone also leaves the
stability question open: the sum cancellation holds even for a time step
that amplifies oscillating modes.
